%%% exercise 04 %%%
\documentclass{article}
\usepackage[margin=1in]{geometry}
\usepackage{amsmath, amssymb, amsfonts, enumitem, fancyhdr, tikz}
% get rid of paragraph indent
\setlength{\parindent}{0 pt}
% allow section.equation numbering
\numberwithin{equation}{section}
% allows you to copy-paste code segments. requires pygments, which can be
% installed from PyPI with pip install Pygments.
% warning: minted does NOT work with Texmaker if you have TeX Live installed
% on WSL but Texmaker installed natively on Windows!
%\usepackage{minted}
% alternative to minted that does not require Python, LaTeX only. listings is
% however disgusting out of the box and some setup is required.
\usepackage{listings, xcolor}
% makes clickable links to sections
\usepackage{hyperref}
% make the link colors blue, as well as cite colors. urls are magenta
\hypersetup{
    colorlinks, linkcolor = blue, citecolor = blue, urlcolor = magenta
}
% fancy pagestyle so we can use fancyhdr for fancy headers/footers
\pagestyle{fancy}
% add logo in right of header. note that you will have to adjust logo path!
\fancyhead[R]{\includegraphics[scale = 0.15]{../bac_logo1.png}}
% don't show anything in the left and center header
\fancyhead[L, C]{}
% give enough space for logo by reducing top margin height, head separator,
% increasing headerheight. see Figure 1 in the fancyhdr documentation. if
% \topmargin + \headheight + \headsep = 0, original text margins unchanged.
\setlength{\topmargin}{-60 pt}
\setlength{\headheight}{50 pt}
\setlength{\headsep}{10 pt}
% remove decorative line in the fancy header
\renewcommand{\headrulewidth}{0 pt}

% color definitions for listings syntax highlighting. uses colors borrowed
% from the VS Code Dark+ and Abyss standard themes.
\definecolor{KwColor}{RGB}{153, 102, 184}     % keyword color
\definecolor{VarColor}{RGB}{86, 156, 214}     % variables/identifier color
\definecolor{StrColor}{RGB}{209, 105, 105}    % string color
\definecolor{CmtColor}{RGB}{106, 153, 85}     % comment color

% general listings configuration for all languages
\lstset{
    % change keyword, identifier, comment, string colors
    keywordstyle = \color{KwColor},
    commentstyle = \color{CmtColor},
    identifierstyle = \color{VarColor},
    stringstyle = \color{StrColor},
    % no spaces in strings
    showstringspaces = false,
    % monospace font by default
    basicstyle = \ttfamily,
    % tabsize 8 by default, this is not the 1960s
    tabsize = 4,
    % add line numbers to the left with gray typewriter font
    numbers = left,
    numberstyle = \color{gray}\ttfamily,
    % change distance from code block from 10 pt to 5 pt
    numbersep = 5 pt
}

% title, author + thanks, date
\title{Exercise 4}
\author{Derek Huang\thanks{NYU Stern 2021, BAC Advanced Team.}}
\date{March 30, 2021\thanks{Original version released on March 11, 2021.}}

% shortcut links. the % characters strip extra spacing.
\newcommand{\pytest}{\href{https://docs.pytest.org/en/stable/}{pytest}}
\newcommand{\minimize}{%
    \href{%
        https://docs.scipy.org/doc/scipy/reference/generated/%
        scipy.optimize.minimize.html%
    }{\texttt{scipy.optimize.minimize}}%
}
% don't wanna keep typing this out
\newcommand{\loglossgrad}{\texttt{\_logistic\_loss\_grad}}

\begin{document}

\maketitle
% need to include this after making title to undo the automatic
% \thispagestyle{plain} command that is issued.
\thispagestyle{fancy}


\section{Introduction}

The goal of this exercise is to implement an $ \ell^2 $-regularized logistic
regression model, minimizing the convex objective using the L-BFGS
implementation provided by \minimize.

\section{Instructions}

\subsection{General}

The \texttt{exercise\_04.py} file contains a skeleton for the
\texttt{LogisticRegression} class, the \loglossgrad{} function, a unit test,
and a \pytest{} fixture. Your job is to implement the \texttt{predict} and
\texttt{score} methods of the \texttt{LogisticRegression} class and the
\loglossgrad{} function by replacing the
\texttt{\#\#\# your code goes here \#\#\#} comment blocks with the appropriate
Python code.

\medskip

Your code \textbf{must} be written in the areas marked by these blocks. Do
\textbf{not} change any of the pre-written code. The exercise is complete
$ \Leftrightarrow $ \texttt{pytest /path/to/exercise\_04.py} executes with
zero test failures. Note that attributes ending with \texttt{\_}, ex.
\texttt{coef\_}, are created \textbf{during} the fitting process.

\subsection{Objective and gradient}

Let $ \mathbf{X} \in \mathbb{R}^{N \times d} $ be the input matrix, each row
$ \mathbf{x}_k^\top $ a tranposed data example, and let $ \mathbf{y} \in
\{-1, 1\}^N $ be the response label vector, $ -1 $ for negative class, $ 1 $
for positive class. Let $ \mathbf{w} \in \mathbb{R}^d $ be the weight vector,
$ b \in \mathbb{R} $ the intercept term. Then, for $ \alpha \in [0, \infty) $,
the loss functional $ J_\mathcal{D} : \mathbb{R}^d \times \mathbb{R}
\rightarrow [0, \infty) $, $ \mathcal{D} \triangleq
(\mathbf{X}, \mathbf{y}) $, is such that
\begin{equation} \label{logreg_obj}
    J_\mathcal{D}(\mathbf{w}, b) = \sum_{k = 1}^N\log\left(
        1 + e^{-y_k(\mathbf{w}^\top\mathbf{x}_k + b)}
    \right) + \frac{1}{2}\alpha\Vert\mathbf{w}\Vert_2^2
\end{equation}

This form is exactly the unweighted objective minimized by scikit-learn for
two-class logistic regression. To minimize $ J_\mathcal{D} $, we will be using
L-BFGS, a \href{%
    https://en.wikipedia.org/wiki/Quasi-Newton_method%
}{quasi-Newton method} that is a limited-memory version of the \href{%
    https://en.wikipedia.org/wiki/Broyden\%E2\%80\%93Fletcher\%E2\%80\%93%
    Goldfarb\%E2\%80\%93Shanno_algorithm%
}{BFGS algorithm}. The $ \mathbf{w}, b $ gradients of $ J_\mathcal{D} $, which
are required\footnote{
    If the gradient is not provided, \minimize{} will approximate it with
    finite differences.
} by [L-]BFGS, are such that
\begin{equation} \label{logreg_grads}
    \begin{split}
        \nabla_\mathbf{w}J_\mathcal{D}(\mathbf{w}, b) & =
	    \sum_{k = 1}^N y_k\mathbf{x}_k\left[
	        \sigma\left(
	            y_k\left(\mathbf{w}^\top\mathbf{x}_k + b\right)
	        \right) - 1
	    \right] + \alpha\mathbf{w} \\
	    \nabla_bJ_\mathcal{D}(\mathbf{w}, b) & =
	    \sum_{k = 1}^N y_k\left[
	        \sigma\left(
	            y_k\left(\mathbf{w}^\top\mathbf{x}_k + b\right)
	        \right) - 1
	    \right]
    \end{split}
\end{equation}

Here $ \sigma : \mathbb{R} \rightarrow (0, 1) $ is the sigmoid function, i.e.
$ \sigma(x) \triangleq 1 / (1 + e^{-x}) $.

\subsection{Implementing \loglossgrad}

\loglossgrad{} computes the loss and gradient for a two-class logistic
regression model. It takes positional parameters \texttt{w}, the concatenated 
coefficients $ [ \ \mathbf{w} \ \ b \ ]^\top \in \mathbb{R}^{d + 1} $,
\texttt{X}, the input matrix $ \mathbf{X} \in \mathbb{R}^{N \times d} $,
\texttt{y}, the label vector $ \mathbf{y} \in \{-1, 1\}^N $, and
\texttt{alpha}, the $ \alpha \in [0, \infty) $ parameter scaling the
$ \ell^2 $ regularization term. Note that \texttt{alpha} does \textbf{not}
show up in the signature of \texttt{LogisticRegression.\_\_init\_\_} and so
\texttt{1 / self.C} should be passed to the \texttt{alpha} parameter of
\loglossgrad, where \texttt{C} is the inverse regularization parameter
mentioned in the class docstring. The use of \texttt{C} here is to conform to
scikit-learn convention.

\medskip

The tuple \texttt{(loss, grad)} should be returned by \loglossgrad{}, where
\texttt{loss} gives the current loss functional value
$ J_\mathcal{D}(\mathbf{w}, b) $ and \texttt{grad} gives the loss functional
$ [ \ \mathbf{w} \ \ b \ ]^\top $ gradient
$ [ \ \nabla_\mathbf{w}J_\mathcal{D}(\mathbf{w}, b)^\top \ \
\nabla_bJ_\mathcal{D}(\mathbf{w}, b)^\top \ ]^\top $. \texttt{loss} must be of
type \texttt{float} or \texttt{numpy.float64} and \texttt{grad} must be of
type \texttt{numpy.ndarray} with a float \texttt{dtype}.

\subsection{Implementing \texttt{predict}}

The \texttt{predict} method of the \texttt{LogisticRegression} class, when
given \texttt{X}, a \texttt{numpy.ndarray} shape
\texttt{(n\_samples, n\_features)}, where \texttt{n\_features} $ = d $,
should return a \texttt{numpy.ndarray} shape \texttt{(n\_samples,)} giving the
predicted class labels, containing only the values in \texttt{self.classes\_}
recorded by the \texttt{fit} method. Given a new example $ \mathbf{x} \in
\mathbb{R}^d $ and parameters $\mathbf{w} \in \mathbb{R}^d $, $ b \in
\mathbb{R} $, the prediction rule $ F : \mathbb{R}^d \rightarrow
\{\gamma_{-1}, \gamma_1\} $ is such that
\begin{equation*}
    F(\mathbf{x}) \triangleq \gamma_1\mathbb{I}_{(0, \infty)}
        \left(\mathbf{w}^\top\mathbf{x} + b\right) +
        \gamma_{-1}\mathbb{I}_{(-\infty, 0]}
        \left(\mathbf{w}^\top\mathbf{x} + b\right)
\end{equation*}

Here $ \gamma_1 $, the positive class label, should be
\texttt{self.classes\_[1]}, while $ \gamma_{-1} $, the negative class label,
should be \texttt{self.classes\_[0]}. Be careful not to get them mixed up or
a correct implementation will predict the opposite of the scikit-learn
implementation for any new data example $ \mathbf{x} $.

\subsection{Implementing \texttt{score}}

The \texttt{score} method of the \texttt{LogisticRegression} class computes
the accuracy of the predictions. Given parameters $ \mathbf{w}, b $,
classification function $ F_{\mathbf{w}, b} $ (the model), and data
$ \mathcal{D} \triangleq \{(\mathbf{x}_1, y_1), \ldots (\mathbf{x}_N, y_N)\} $,
where the $ y_k $ values may either be numeric or class labels, the accuracy
$ \mathcal{A}_{\mathbf{w}, b} $ of the model on the data $ \mathcal{D} $ is
\begin{equation*}
    \mathcal{A}_{\mathbf{w}, b}(\mathcal{D}) = \frac{1}{|\mathcal{D}|}
        \sum_{(\mathbf{x}, y) \in \mathcal{D}}\mathbb{I}_{\{y\}}\left(
            F_{\mathbf{w}, b}(\mathbf{x})
        \right)
\end{equation*}

Note that $ \operatorname{im}\mathcal{A}_{\mathbf{w}, b} = [0, 1] $, so for
data $ \mathcal{D} $, $ \mathcal{A}_{\mathbf{w}, b}(\mathcal{D}) $ is the
fraction of examples the model correctly clasified.

\subsection{Tips}

\begin{enumerate}
    \item
    Carefully read the \texttt{LogisticRegression} method docstrings.

    \item
    Remember that \loglossgrad{} should return both the objective (loss
    functional) value \textbf{and} its $ [ \ \mathbf{w} \ \ b \ ]^\top $
    gradient. If \texttt{n\_features} $ = d $, allocate a
    \texttt{numpy.ndarray} shape \texttt{(n\_features + 1,)}, filling the
    first \texttt{n\_features} elements with the $ \mathbf{w} $ gradient
    and the last with the $ b $ gradient (partial derivative).

    \item
    \href{%
        https://docs.scipy.org/doc/scipy/reference/generated/%
        scipy.special.expit.html%
    }{\texttt{scipy.special.expit}} is a vectorized elementwise sigmoid
    function that is very useful for computing the $ \mathbf{w}, b $
    gradients defined in (\ref{logreg_grads}) for \loglossgrad.

    \item
    Remember that $ \alpha = 1 / C $, so $ 1 / C $ should be passed to the
    \texttt{alpha} parameter of \loglossgrad.

    \item
    Use NumPy functions whenever possible for efficiency, brevity, and ease
    of debugging.
\end{enumerate}




\end{document}